\section{Methodology}
\label{sec:method}

\subsection{Problem formulation}

Let $\vp = (p_{\text{web}}, p_{\text{wiki}}, p_{\text{code}}, p_{\text{arxiv}})$ be a mixture on the simplex, where $p_d$ is the fraction of training tokens drawn from domain $d$.
We fix the architecture, tokenizer, token budget, optimizer, and schedule, and train one model per mixture.
We then ask how a score $y(\vp)$, either held-out loss on a domain or accuracy on a downstream task, depends on $\vp$.

\para{Curve shapes along a sweep.}
For a single-domain sweep we compare four forms for $y$ as a function of the swept share $p$:
\begin{align}
\text{linear:} \quad & y = a + b\,p, \\
\text{log:} \quad & y = a + b \log p, \label{eq:log} \\
\text{exponential:} \quad & y = c + k\,e^{t p}, \\
\text{log with plateau:} \quad & y = a + b \log \min(p, p^{\star}).
\end{align}
The exponential form is the one-dimensional version of the law of \citet{ye2025datamixing}, and $p^{\star}$ is a fitted saturation point.
H1 says that the log form beats the linear form.
H2 says that $p^{\star} \approx 0.2$.

\para{Laws over the whole simplex.}
To predict unseen mixtures we compare a linear (Scheff\'e) mixture model $y = \sum_j \beta_j p_j$, the log law in the domain's own share (\eqref{eq:log}), the multi-domain exponential law $y = c + k \exp(\sum_j t_j p_j)$ of \citet{ye2025datamixing}, and an additive log-share law
\begin{equation}
L_d(\vp) = a_d + \sum_{j} b_{dj} \log p_j ,
\label{eq:additive}
\end{equation}
where $L_d$ is held-out loss on domain $d$.
The coefficients $b_{dj}$ form a transfer matrix: $b_{dj}$ is the change in loss on domain $d$ per $e$-fold increase in the share of domain $j$.

\subsection{Pre-training}

\para{Domains.}
We use four domains: \web is English C4 \citep{raffel2020t5}; \wiki is WikiText-103 \citep{merity2017wikitext}; \code is the Python corpus \texttt{codeparrot-clean}; \arxiv is the article bodies of the arXiv summarization dataset \citep{cohan2018discourse}.
We equalize the pools at 128.9M training tokens per domain and hold out separate documents for evaluation.
One byte-level BPE tokenizer with 16{,}384 tokens, trained on a balanced sample, is shared by all runs.

\para{Model and budget.}
Every model is a decoder-only transformer with 6 layers, width 384, 6 heads, context length 256, and tied embeddings, for 17.0M parameters.
Every run sees exactly 60M tokens (3{,}662 steps of 64 sequences of 256 tokens).
We train with AdamW \citep{loshchilov2019adamw} at learning rate $10^{-3}$, 3\% warm-up, cosine decay to 10\% of the peak, weight decay 0.1, $\beta = (0.9, 0.95)$, gradient clipping at 1.0, and bfloat16 autocast.
Each training sequence is drawn from domain $d$ with probability $p_d$.
We chose this size because the GPUs were shared with other jobs; \secref{sec:limitations} and Appendix~\ref{app:details} describe this and other deviations from the original plan.

\subsection{Mixture design}

We train 81 runs on 53 distinct mixtures (\tabref{tab:design}).

\begin{table}[t]
\centering
\caption{Mixture design. The uniform mixture lies on all four sweeps. Seed replicates use three seeds each and give 28 degrees of freedom for the noise estimate.}
\label{tab:design}
\small
\begin{tabular}{@{}p{0.17\textwidth}p{0.45\textwidth}p{0.31\textwidth}@{}}
\toprule
Component & Mixtures & Purpose \\
\midrule
Single-domain sweeps & share $p \in \{0, 0.02, 0.05, 0.1, 0.2,$ $0.35, 0.5, 0.75, 1\}$ per domain; the other three split $1-p$ equally & curve shape, saturation, web threshold (H1, H2, H4) \\
Random mixtures & 16 draws from $\mathrm{Dirichlet}(1,1,1,1)$ & unbiased correlations (H3); held-out test set for laws \\
Seed replicates & uniform; every $p{=}0$ and $p{=}1$ model; \web at 0.2, 0.5, 0.75; \code at 0.2, 0.5 (14 mixtures) & noise floor \\
\bottomrule
\end{tabular}
\end{table}

\para{Why three components?}
The sweeps are the only way to see the shape of a curve and to locate a saturation point.
But shares on a simplex are negatively correlated with one another, and a sweep confounds a domain's effect with the sweep structure, so we use Dirichlet draws for correlations.
Without seed replicates, no claim about curve shape or degradation is interpretable.

\subsection{Evaluation}

\para{Loss.}
We compute held-out loss in nats per token on each domain over non-overlapping windows (0.27--1.0M tokens per domain).

\para{Knowledge, zero-shot ($\Kzs$).}
We score multiple-choice questions by choosing the answer with the highest length-normalized log-likelihood after the prompt ``\texttt{Question:\,\ldots\textbackslash nAnswer:}''.
$\Kzs$ is the mean accuracy on \sciq \citep{welbl2017sciq} (13{,}679 items) and \arceasy \citep{clark2018arc} (5{,}197 items).
We also report ARC-Challenge, OpenBookQA \citep{mihaylov2018openbookqa}, and CommonsenseQA \citep{talmor2019commonsenseqa} but exclude them from the composite.

\para{Knowledge, fine-tuned ($\Kft$).}
We fine-tune each checkpoint for 400 steps on \sciq training answers without the supporting passage, then score multiple choice on the 2{,}000 validation and test items.

\para{Reasoning ($\Rc$).}
Models of this size cannot reason zero-shot, so we measure fine-tuning transfer.
$\Rc$ is the mean exact-match accuracy on \listops \citep{nangia2018listops} (generated; 4{,}000 training and 2{,}000 test examples) and \proofwriter \citep{tafjord2021proofwriter} (3{,}000 training and 1{,}500 test examples), each after 500 fine-tuning steps.
We also report 3-digit arithmetic after 1{,}500 steps, outside the composite.

Fine-tuning uses learning rate $3 \times 10^{-4}$, loss on answer tokens only, and the same data order for every checkpoint.
Re-running the evaluation of a checkpoint reproduces every number exactly.

\para{How we chose the composites.}
We fixed task inclusion using only the uniform-mixture pilot and a randomly initialized control, before any mixture comparison.
$\Kzs$ keeps the two zero-shot tasks that were above chance.
$\Kft$ is the fallback we planned in case zero-shot signal was weak.
$\Rc$ keeps the two fine-tuned tasks on which pre-training helped.
We exclude arithmetic because pre-training hurt it in the pilot.

\subsection{Baselines and statistical analysis}

\para{Baselines.}
We compare against the uniform mixture, pure-domain models ($p = 1$), leave-one-domain-out models ($p = 0$), a randomly initialized model for all fine-tuned probes, and chance for multiple choice.
For curve shape the linear fit is the null.

\para{Tests.}
For curve shape we use leave-one-out (LOO) RMSE as the primary criterion and AICc as a check, fitting on the nine shares above zero.
For correlations we report Pearson $r$ with bootstrap 95\% confidence intervals and a Holm correction across the four domains.
For degradation we use Welch $t$-tests across seeds.
The noise floor is the pooled standard deviation across seeds at the 14 replicated mixtures.
We fit simplex laws on the 29 structured mixtures in which every share is positive and test on the 16 random mixtures.
